% GENERATED FILE. Edit canonical lessons, then run npm run generate:books.
% canonical-source-hash: fnv1a64:19aff8dc86b2bdd0
% canonical-lessons: IT-C437-baccano, IT-C437-chiasso, IT-C437-scarafaggio, IT-C437-muffa, IT-C437-cigolare

\chapter{Rumore e sporco}
\label{ch:it-rumore-e-sporco}
\hlchaptermodality{437}

\begin{chapteropening}
\textbf{By the end of this chapter:} \emph{I can complain about noise and dirt: a racket, a din, a cockroach, mould and a creaking door.}

Complete the last lesson of chapter 437: I can complain about noise and dirt: a racket, a din, a cockroach, mould and a creaking door.
\end{chapteropening}

\section[il baccano]{il baccano — the racket, the uproar}

\label{lesson:IT-C437-baccano}



\begin{quote}
Before the new one: say the Italian for to recharge, then the Italian for to burn out.
\end{quote}

\subsection*{You'll want to know: il baccano}
\textbf{il baccano} — \textquotedblleft{}the racket, the uproar\textquotedblright{}.

\begin{quote}
I vicini fanno un baccano tremendo.
\end{quote}

\begin{quote}
\emph{The neighbours are making a dreadful racket.}
\end{quote}

Informal; \emph{fare baccano} — to make a racket.

\begin{cousinweb}
\emph{baccano} is built on \emph{baccanale} 'a wild party', a learned borrowing from Latin \emph{bacchānal}, a feast of the Greek wine god (EtymDB).
\end{cousinweb}

\begin{grammarlens}[title={What you've built}]
One more word for telling, explaining and describing at B1.
\end{grammarlens}

\subsection*{Your turn}
\begin{itemize}
\raggedright
  \item \emph{Say it:} \emph{il baccano}
  \item \emph{Say it:} \emph{il baccano}, once more
  \item \emph{Say it:} a sentence of your own with \emph{il baccano}
  \item \emph{Recall:} say the Italian for to recharge, then the Italian for to burn out, then say \emph{il baccano} again
\end{itemize}

\begin{tcolorbox}[breakable,colback=teal!4,colframe=teal!35!black,title={Before you move on}]
What does il baccano mean? (\textquotedblleft{}The racket, the uproar\textquotedblright{}.) Say it once more.
\end{tcolorbox}
\section[il chiasso]{il chiasso — the din, the noise (of voices)}

\label{lesson:IT-C437-chiasso}



\begin{quote}
Before the new one: say the Italian for to burn out, then the Italian for the racket.
\end{quote}

\subsection*{You'll want to know: il chiasso}
\textbf{il chiasso} — \textquotedblleft{}the din, the noise (of voices)\textquotedblright{}.

\begin{quote}
Al mercato c'è sempre chiasso.
\end{quote}

\begin{quote}
\emph{At the market there is always a din.}
\end{quote}

\emph{Fare chiasso} — to make a noise; \emph{chiassoso} — noisy.

\begin{cousinweb}
\emph{chiasso} is usually traced to Latin \emph{classicum} 'trumpet call' (Lewis \& Short), but the origin is not settled.
\end{cousinweb}

\begin{grammarlens}[title={What you've built}]
One more word for telling, explaining and describing at B1.
\end{grammarlens}

\subsection*{Your turn}
\begin{itemize}
\raggedright
  \item \emph{Say it:} \emph{il chiasso}
  \item \emph{Say it:} \emph{il chiasso}, once more
  \item \emph{Say it:} a sentence of your own with \emph{il chiasso}
  \item \emph{Recall:} say the Italian for to burn out, then the Italian for the racket, then say \emph{il chiasso} again
\end{itemize}

\begin{tcolorbox}[breakable,colback=teal!4,colframe=teal!35!black,title={Before you move on}]
What does il chiasso mean? (\textquotedblleft{}The din, the noise (of voices)\textquotedblright{}.) Say it once more.
\end{tcolorbox}
\section[lo scarafaggio]{lo scarafaggio — the cockroach}

\label{lesson:IT-C437-scarafaggio}



\begin{quote}
Before the new one: say the Italian for the racket, then the Italian for the din.
\end{quote}

\subsection*{You'll want to know: lo scarafaggio}
\textbf{lo scarafaggio} — \textquotedblleft{}the cockroach\textquotedblright{}.

\begin{quote}
In cucina c'era uno scarafaggio!
\end{quote}

\begin{quote}
\emph{There was a cockroach in the kitchen!}
\end{quote}

A good reason to ask for another room.

\begin{cousinweb}
\emph{scarafaggio} is usually traced to Latin \emph{scarabaeus} 'beetle' (Lewis \& Short), from Greek \emph{kárabos}; it is inherited, with a changed ending.
\end{cousinweb}

\begin{grammarlens}[title={What you've built}]
One more word for telling, explaining and describing at B1.
\end{grammarlens}

\subsection*{Your turn}
\begin{itemize}
\raggedright
  \item \emph{Say it:} \emph{lo scarafaggio}
  \item \emph{Say it:} \emph{lo scarafaggio}, once more
  \item \emph{Say it:} a sentence of your own with \emph{lo scarafaggio}
  \item \emph{Recall:} say the Italian for the racket, then the Italian for the din, then say \emph{lo scarafaggio} again
\end{itemize}

\begin{tcolorbox}[breakable,colback=teal!4,colframe=teal!35!black,title={Before you move on}]
What does lo scarafaggio mean? (\textquotedblleft{}The cockroach\textquotedblright{}.) Say it once more.
\end{tcolorbox}
\section[la muffa]{la muffa — the mould}

\label{lesson:IT-C437-muffa}



\begin{quote}
Before the new one: say the Italian for the din, then the Italian for the cockroach.
\end{quote}

\subsection*{You'll want to know: la muffa}
\textbf{la muffa} — \textquotedblleft{}the mould\textquotedblright{}.

\begin{quote}
Sul muro del bagno c'è la muffa.
\end{quote}

\begin{quote}
\emph{There is mould on the bathroom wall.}
\end{quote}

\emph{Sapere di muffa} — to smell musty.

\begin{cousinweb}
\emph{muffa} is usually traced to a Germanic word for mould or a musty smell, but the origin is not settled.
\end{cousinweb}

\begin{grammarlens}[title={What you've built}]
One more word for telling, explaining and describing at B1.
\end{grammarlens}

\subsection*{Your turn}
\begin{itemize}
\raggedright
  \item \emph{Say it:} \emph{la muffa}
  \item \emph{Say it:} \emph{la muffa}, once more
  \item \emph{Say it:} a sentence of your own with \emph{la muffa}
  \item \emph{Recall:} say the Italian for the din, then the Italian for the cockroach, then say \emph{la muffa} again
\end{itemize}

\begin{tcolorbox}[breakable,colback=teal!4,colframe=teal!35!black,title={Before you move on}]
What does la muffa mean? (\textquotedblleft{}The mould\textquotedblright{}.) Say it once more.
\end{tcolorbox}
\section[cigolare]{cigolare — to creak, to squeak}

\label{lesson:IT-C437-cigolare}



\begin{quote}
Before the new one: say the Italian for the cockroach, then the Italian for the mould.
\end{quote}

\subsection*{You'll want to know: cigolare}
\textbf{cigolare} — \textquotedblleft{}to creak, to squeak\textquotedblright{}.

\begin{quote}
La porta cigola ogni volta che la apri.
\end{quote}

\begin{quote}
\emph{The door creaks every time you open it.}
\end{quote}

\emph{Il cigolio} — creaking.

\begin{cousinweb}
\emph{cigolare} seems to imitate the squeak itself; the origin is not settled.
\end{cousinweb}

\begin{grammarlens}[title={What you've built}]
One more word for telling, explaining and describing at B1.
\end{grammarlens}

\subsection*{Your turn}
\begin{itemize}
\raggedright
  \item \emph{Say it:} \emph{cigolare}
  \item \emph{Say it:} \emph{cigolare}, once more
  \item \emph{Say it:} a sentence of your own with \emph{cigolare}
  \item \emph{Recall:} say the Italian for the cockroach, then the Italian for the mould, then say \emph{cigolare} again
\end{itemize}

\begin{tcolorbox}[breakable,colback=teal!4,colframe=teal!35!black,title={Before you move on}]
What does cigolare mean? (\textquotedblleft{}To creak, to squeak\textquotedblright{}.) Say it once more.
\end{tcolorbox}
